\documentclass[11pt]{beamer}

% \setbeamersize{text margin left=5mm, text margin right=5mm}

%%%%%%%%%%%%%%%%%%%%%%%%%
%   Options for notes   %
%%%%%%%%%%%%%%%%%%%%%%%%%
% \setbeamertemplate{note page}[plain]
% \setbeameroption{show notes}
% \setbeameroption{show notes on second screen}
%\setbeameroption{show only notes}
% \setbeameroption{show only slides with notes}

\input{preamble}

%%%%%%%%%%%%%%%%%%%%%%%%%
%     presenter info    %
%%%%%%%%%%%%%%%%%%%%%%%%%
\author[Xinda GE]{
Xinda GE\\
\small PhD year 1 \\ Mentor: Stephen Dalton
}
\title[SOX2 in hESC by Degron system]{
How SOX2 as a pioneer factor to initiate ESCs differentiation}

% \subtitle{
% Mentor: Dalton
% }
\institute[Dalton Lab]{
Dalton Lab,\\
School of Biomedical Sciences
}
\date{\advanceday[1]}


\begin{document}

% \kaishu
\begin{frame}
    \titlepage
    \begin{figure}[htpb]
        \begin{center}
            
\includegraphics[width=0.1\linewidth]{pic/CUHKSZ-Logo.png}
        \end{center}
    \end{figure}
\end{frame}

%%%%%%%%%%%%%%%%%%% Contents Start %%%%%%%%%%%%%%%%%%%%%%
\section{Background}

\begin{frame}{Why SOX2 might be essential for ESCs differentiation}  

    \begin{figure}
        \centering
        \includegraphics[width=1.0\textwidth]{20230422-SOX2/figs/background.pdf}
        %caption{Approaches to clear protein molecules inside the nucleus of a cell}
        %label{fig:ubprocess}
    \end{figure}
    
    \begin{itemize}
        \item SOX2 is engaged with the three germ layers' differentiation.
    \end{itemize}
    
\end{frame}


% \note[enumerate]{
%     \item ICM cells derived from the Blastoid has several potentials to differentiate into 3 germ layers, en,me, ec. 
%     \item Of notes, OCT4, SOX2 and Nanog as 3 key transcription factors are essential for the pluripotency state.
%     \item Here is an image to show the Chip-seq result which reflects the TF binding frequency onto the target motif. 
%     \item 
% }

\section{Method}

\begin{frame}{Inducible KO strategies of a targeted protein in nucleus}  
    % \begin{itemize}
    %     \item 
    % \end{itemize}
    \begin{columns}
        \column{0.6\textwidth}
        \begin{figure}
            \centering
            \includegraphics[width=\textwidth]{20230422-SOX2/figs/xinda.drawio.pdf}
            %caption{Approaches to clear protein molecules inside the nucleus of a cell}
            %label{fig:ubprocess}
        \end{figure}
        \column{0.4\textwidth}
        
        Advantages of degron system
        % {\small
        \begin{itemize}
            \item without disrupting pl-uripotency
            \item accurately on-target
            \item reversible
            \item rapid 
            \item $\cdots$
        \end{itemize}
        % }
    \end{columns}
\end{frame}




\begin{frame}{What are needed to create a Degron cell line?}  
    % \begin{itemize}
    %     \item 
    % \end{itemize}
        
    \begin{figure}
        \centering
        \includegraphics[width=\textwidth]{20230422-SOX2/figs/howdegron.pdf}
        %caption{AID system mechanism}
        %label{fig:mechanisms}
    \end{figure}

    \begin{itemize}
        \item The key components include \textbf{TIR1}, \textbf{AID tag}, and \textbf{auxin}.
        \item The advantages of degron system are confirmed.
    \end{itemize}
\end{frame}




\section{Experiment Design}

\subsection{AID-tag onto SOX2 + Validation}

\begin{frame}{1. Tag SOX2 with AID of Degron system by CRISPR/Cas9}  
    % \begin{itemize}
    %     \item  
    % \end{itemize}
        
    \begin{figure}
        \centering
        \includegraphics[width=0.9\textwidth]{20230422-SOX2/figs/design1.pdf}
        %caption{Gene Editing Process}
        %label{fig:gene editing}
    \end{figure}
\end{frame}

\begin{frame}{Insertion validation (Part 1)}  
    % \begin{itemize}
    %     \item  
    % \end{itemize}
        
    \begin{figure}
        \centering
        \includegraphics[width=\textwidth]{20230422-SOX2/insertvalp1.png}
        %caption{Gene insertion evidence}
        %label{fig:gene insertion}
    \end{figure}
\end{frame}

\begin{frame}{Insertion validation (Part 2)}  
    % \begin{itemize}
    %     \item  
    % \end{itemize}
        
    \begin{figure}
        \centering
        \includegraphics[width=1.0\textwidth]{20230422-SOX2/figs/insertvalp2.pdf}
        %caption{Protein expression evidence}
        %label{fig:gene editing}
    \end{figure}
\end{frame}

\subsection{TIR1 into AAVS1 + Validation}

\begin{frame}{2. Safe harbor locus insertion with TIR1}  
    % \begin{itemize}
    %     \item  
    % \end{itemize}
        
    \begin{figure}
        \centering
        \includegraphics[width=1.0\textwidth]{20230422-SOX2/figs/safeharbor.pdf}
        %\caption{Gene Editing Process}
        %label{TIR1 insertion into AAVS1 locus}
    \end{figure}
\end{frame}

\begin{frame}{Proof-of-concept validation}  
    \begin{figure}
        \centering
        \includegraphics[width=1.0\textwidth]{20230422-SOX2/figs/proofval.png}
        %\caption{Gene Editing Process}
        \label{Function validation}
        \href{/Users/gexinda/Downloads/OneDrive_1_7-30-2023/500nM.avi}{500 nM IAA.avi}\\
        \href{/Users/gexinda/Downloads/OneDrive_1_7-30-2023/1uM.avi}{1 uM IAA.avi}
    \end{figure}
    \begin{itemize}
        \item Transient transfection Degron2 (2nd generation Degron)
        (merged images)
    \end{itemize}
\end{frame}

\section{Summary}
\begin{frame}{Summary} 
 \begin{enumerate}
       \item SOX2-AID-GFP homozygous cell line is generated;
       \item SOX2 can be degraded upon Degron2 system;
       \item We are pending on Degron2 KI to test the function;
       \item This cell line will be a flexible tools to investigate the function of SOX2 in ESCs differentiation.
       \end{enumerate}
\end{frame}



%%%%%%%%%%%%%%%%%%% Contents End %%%%%%%%%%%%%%%%%%%%%%

%
% section*{References}
% \begin{frame}{References}
%     \footnotesize
%     \nocite{*}
%     \bibliography{ref}
%     \bibliographystyle{abbrv}
%     % \tiny\bibliographystyle{alpha}
% \end{frame}


\begin{frame}
    \begin{center}
        {\huge Acknowledgements}
        \begin{align*}
            \text{Prof.}~& \text{Stephen DALTON,} \\
            \text{Dr.}~& \text{Tianming WU,} \\\\
            \text{Dr.}~& \text{Dong LI,} \\
            \text{Ms.}~& \text{Sau Kuen NG,} \\
            \text{Dr.}~& \text{Angel Yuet Fong KAM,} \\
            \text{Mr.}~& \text{Hao YU,} \\
            \text{Miss}~& \text{Kathy SHAM,} \\
            \text{Mr.}~& \text{Xiaocheng ZHOU}
        \end{align*}
    \end{center}
\end{frame}

\end{document}